% 全书的版式（仿照《Computer Organization and Design》第 5 版）：这一部分放在导言区。
% 真正启用版式的是 style/lazy-on.tex。
%   章首页：左上角一个红色方块，里面是白色的大号章号；一条竖线；
%           右边是黑体的章标题和本章各节的小目录，左下是一段引言
%   页眉：  偶数页“页码  第一章  章标题”，奇数页“节标题  页码”，下面一条红线
%   页脚：  中文版偶数页英文书名、奇数页中文书名；英文版偶数页书名、奇数页作者
%   节标题：红底白字的编号方块，后面是黑体标题
\usepackage{fancyhdr}
\usepackage{etoc}
\usepackage{mparhack}   % 修正页边注偶尔放错边的问题（段落跨页时，LaTeX 会把注放到内侧）

\definecolor{lazyred}{RGB}{190,30,45}          % 和封面的红色一样
\newfontfamily\lazysans{texgyreheros}[
  Extension=.otf, Ligatures=TeX, UprightFont=*-regular, BoldFont=*-bold,
  ItalicFont=*-italic, BoldItalicFont=*-bolditalic]
\setCJKfamilyfont{lazyhei}{FandolHei-Regular}[Extension=.otf, BoldFont=FandolHei-Bold]
\setCJKfamilyfont{lazykai}{FandolKai-Regular}[Extension=.otf]
\newcommand\lazyfont{\lazysans\CJKfamily{lazyhei}}   % 标题用的无衬线字体：英文 Heros，中文黑体

% 每章章首页的引言：\lazyquote{章号}{中文}{原文}{出处}，附录的章号是 A、B……
% 各章的引言写在语言文件里（style/lang-zh.tex、style/lang-en.tex）；英文版的“中文”一项留空，章首页就只印原文。
\newcommand\lazyquote[4]{%
  \expandafter\def\csname lazyquoteA#1\endcsname{#2}%
  \expandafter\def\csname lazyquoteB#1\endcsname{#3}%
  \expandafter\def\csname lazyquoteC#1\endcsname{#4}}

% 页眉、页脚。页脚的文字由语言文件定义：\lazyfooteven、\lazyfooterodd（正文页），\lazyfooteropener（章首页）
\fancypagestyle{lazy}{%
  \fancyhf{}%
  \fancyhead[LE]{\lazyfont\bfseries\thepage\hspace{1.6em}\leftmark}%
  \fancyhead[RO]{\lazyfont\bfseries\rightmark\hspace{1.6em}\thepage}%
  \fancyfoot[LE]{\lazyfont\footnotesize\color{gray}\lazyfooteven}%
  \fancyfoot[RO]{\lazyfont\footnotesize\color{gray}\lazyfooterodd}%
  \renewcommand\headrule{{\color{lazyred}\hrule height 0.8pt}\vskip-0.8pt}%
  \renewcommand\footrulewidth{0pt}}
% 章首页：没有页眉，页脚居中写书名
\fancypagestyle{lazyopener}{%
  \fancyhf{}%
  \fancyfoot[C]{\lazyfont\footnotesize\color{gray}\lazyfooteropener}%
  \renewcommand\headrule{}}

% 页边注：书页外侧留出一栏空白，放术语解释和小提示（仿照 COD 的做法）
%   \lazynote{术语}{解释}    红色黑体的术语，后面跟一句解释
%   \lazyhint{提示}          灰色的小提示，前面有一个红色的 ▶
\newcommand\lazynotestyle{\lazyfont\footnotesize\linespread{1.15}\selectfont\raggedright}
\newcommand\lazynote[2]{%
  \marginpar{\lazynotestyle{\color{lazyred}\bfseries #1}\quad #2\par}}
\newcommand\lazyhint[1]{%
  \marginpar{\lazynotestyle\color{black!60}{\color{lazyred}$\blacktriangleright$}\ #1\par}}

% 宽的图和表可以伸进页边注那一栏：\begin{lazywide} ... \end{lazywide}
% 放在 figure、table 里，代替 \centering。总宽 = 版心 + 页边注的间隔和宽度。
% 页边注在奇数页的右边、偶数页的左边，所以偶数页要往左挪；\ifoddpage 由 changepage 宏包判断。
\usepackage[strict]{changepage}
\newenvironment{lazywide}
  {\checkoddpage
   \ifoddpage\def\lazywideL{0pt}\def\lazywideR{-50mm}%
   \else     \def\lazywideL{-50mm}\def\lazywideR{0pt}\fi
   \begin{adjustwidth}{\lazywideL}{\lazywideR}\centering
   \captionsetup{width=\linewidth}}%   题注也和加宽后的图一样宽
  {\end{adjustwidth}}

% 节标题的编号：红底白字的方块
\newcommand\lazysecnum[1]{\colorbox{lazyred}{\makebox[2.4em]{\color{white}#1}}}
\newcommand\lazygobble[1]{}

% 章标题：有编号的章（正文各章、附录）用整页的章首页；
% 没有编号的章（前言、目录、代码清单、参考文献）只在页面上方放一个小一些的标题
\newcommand\lazychapter[1]{\CTEXifname{\lazychapternumbered{#1}}{\lazychapterplain{#1}}}

\newcommand\lazychapterplain[1]{%
  \vspace*{6mm}%
  \noindent{\color{lazyred}\rule[-3mm]{5mm}{13mm}}\hspace{5mm}%
  {\lazyfont\bfseries\fontsize{26}{30}\selectfont #1}\par
  \vspace{3mm}%
  \noindent{\color{lazyred}\rule{\linewidth}{0.8pt}}\par
  \vspace{10mm}}

\newcommand\lazychapternumbered[1]{%
  \begin{tikzpicture}[remember picture, overlay]
    \coordinate (P) at (current page.north west);
    % 红色方块和白色的大号章号
    \fill[lazyred] ($(P) + (30mm,-28mm)$) rectangle ($(P) + (112mm,-118mm)$);
    \node[anchor=south west, inner sep=0pt, text=white,
          font=\lazysans\bfseries\fontsize{160}{160}\selectfont]
      at ($(P) + (38mm,-112mm)$) {\thechapter};
    % 竖线
    \draw[lazyred, line width=0.6pt] ($(P) + (112mm,-96mm)$) -- ($(P) + (112mm,-238mm)$);
    % 章标题
    \node[anchor=north west, inner sep=0pt, text width=70mm, align=left,
          font=\lazyfont\bfseries\fontsize{24}{30}\selectfont]
      at ($(P) + (117mm,-124mm)$) (title) {#1};
    % 本章各节的小目录
    \node[anchor=north west, inner sep=0pt, text width=70mm, align=left]
      at ($(title.south west) + (0,-6mm)$) {\lazyminitoc};
    % 引言
    \ifcsname lazyquoteB\thechapter\endcsname
      \node[anchor=north east, inner sep=0pt, text width=68mm, align=flush right]
        at ($(P) + (106mm,-168mm)$)
        {\expandafter\ifx\csname lazyquoteA\thechapter\endcsname\empty
           {\itshape\large\csname lazyquoteB\thechapter\endcsname\par}%   只有原文：原文用大字
         \else
           {\CJKfamily{lazykai}\large\csname lazyquoteA\thechapter\endcsname\par}
           \vspace{2mm}
           {\itshape\small\csname lazyquoteB\thechapter\endcsname\par}%
         \fi
         \vspace{3mm}
         {\lazyfont\bfseries\scriptsize\csname lazyquoteC\thechapter\endcsname\par}};
    \fi
  \end{tikzpicture}%
  \par\vskip\textheight}       % 占满这一页，正文从下一页开始

% 小目录：只列出各节，编号是红色的
\newcommand\lazyminitoc{%
  \begingroup
    \etocsetnexttocdepth{section}%
    \etocsettocstyle{}{}%
    \etocsetstyle{section}{}%
      {}%
      {{\lazyfont\bfseries\small\raggedright\hangindent=2.4em\hangafter=1
        \noindent\makebox[2.4em][l]{\color{lazyred}\etocnumber}\etocname
        \nobreak\hspace{0.6em}{\mdseries\color{gray}\etocpage}\par}\vspace{1.2mm}}%   % 标题太长时折行，悬挂缩进
      {}%
    \localtableofcontents
  \endgroup}
